% TL;DR: Projektkarte als drei Präsentationsfolien 16:9 (Gesamtkarte, Stand heute, Irrwege + Lessons Learned); braucht build/karte_kurz.pdf aus build.sh.
\documentclass{article}
\usepackage[paperwidth=338.67mm, paperheight=190.5mm, margin=0mm]{geometry}
\usepackage{libertinus-otf}
\usepackage{xcolor}
\usepackage{graphicx}
\usepackage{tikz}
% TL;DR: Gemeinsame Farben, TikZ-Stile und Stations-Makros der Projektkarte (U-Bahn-Plan); wird von karte.tex und den Varianten eingebunden.
%
% Lesart der Karte
%   Linie       = Arbeitsstrang (Farbe)
%   Station     = Meilenstein; die Füllung zeigt den Ausgang (\Mk)
%   Stumpfgleis = aufgegebener Weg (endet mit Prellbock)
%   Fehlannahme = gestrichelte Nebenspur (Stil irrweg): so lange galt eine falsche Annahme

\usetikzlibrary{positioning, calc, patterns, shapes.geometric, fit, backgrounds, arrows.meta}

% ── Linienfarben (Arbeitsstränge) ─────────────────────────────────────
\definecolor{lInfra}{HTML}{6B7B8C}   % Infrastruktur / Rechenzugang
\definecolor{lTrain}{HTML}{1F4E9C}   % Training N1.6 (NavyBlue wie die Ausarbeitung)
\definecolor{lNsieben}{HTML}{8E44AD} % GR00T N1.7
\definecolor{lCotr}{HTML}{D1477A}    % Co-Training / Daten
\definecolor{lSim}{HTML}{0E8C7F}     % Simulation / Szene
\definecolor{lDiag}{HTML}{E07B00}    % Evaluation / Diagnose
\definecolor{lRL}{HTML}{8D8D8D}      % Reinforcement Learning
\definecolor{lDoku}{HTML}{3C3C3C}    % Doku / Abgabe

% ── Statusfarben (Stationen) ──────────────────────────────────────────
\definecolor{sOk}{HTML}{2E8B3E}      % Erfolg
\definecolor{sFail}{HTML}{C62828}    % Fehlschlag
\definecolor{sPart}{HTML}{F2B705}    % Teilerfolg
\definecolor{sOld}{HTML}{9E9E9E}     % verworfen / überholt
\definecolor{phaseA}{HTML}{F4F6F9}
\definecolor{phaseB}{HTML}{E9EDF3}
\definecolor{hubbg}{HTML}{FFF8E1}

% ── Maße (die Varianten skalieren die fertige Karte als Ganzes) ───────
\newcommand\linew{2.6pt}
\newcommand\fsLbl{\fontsize{6.4}{7.3}\selectfont\sffamily}
\newcommand\fsPhase{\fontsize{7.2}{8.4}\selectfont\sffamily}

\tikzset{
  line/.style   = {line width=\linew, line cap=round, line join=round, rounded corners=4mm},
  stub/.style   = {line width=\linew, line cap=butt, dash pattern=on 3pt off 2pt},
  irrweg/.style = {line width=1.1pt, dash pattern=on 2.4pt off 1.6pt, line cap=round,
                   rounded corners=2.5mm},
  lbl/.style    = {font=\fsLbl, align=center, inner sep=0.6pt, text depth=0pt},
  badge/.style  = {rounded corners=1.6mm, font=\fontsize{6.8}{7.6}\selectfont\sffamily\bfseries,
                   text=white, inner xsep=1.6mm, inner ysep=1.1mm, minimum width=19mm,
                   align=center},
}

% \Mk{status}{x,y}: Stationsmarke. status ∈ ok | fail | part | old | open
\newcommand\Mk[2]{%
  \begin{scope}[shift={(#2)}]
    \fill[white] (0,0) circle (1.95mm);
    \csname mk@#1\endcsname
  \end{scope}}
\makeatletter
\newcommand\mk@ok{\fill[sOk] (0,0) circle (1.45mm);
  \draw[white, line width=0.85pt, line cap=round, line join=round]
    (-0.72mm,0.02mm) -- (-0.18mm,-0.58mm) -- (0.78mm,0.62mm);}
\newcommand\mk@fail{\fill[sFail] (0,0) circle (1.45mm);
  \draw[white, line width=0.85pt, line cap=round]
    (-0.6mm,-0.6mm) -- (0.6mm,0.6mm) (-0.6mm,0.6mm) -- (0.6mm,-0.6mm);}
\newcommand\mk@part{\fill[white] (0,0) circle (1.45mm);
  \fill[sPart] (0,-1.45mm) arc (-90:90:1.45mm) -- cycle;
  \draw[sPart!80!black, line width=0.6pt] (0,0) circle (1.2mm);}
\newcommand\mk@old{\fill[sOld] (0,0) circle (1.45mm);
  \draw[white, line width=0.8pt] (-0.75mm,0) -- (0.75mm,0);}
\newcommand\mk@open{\fill[white] (0,0) circle (1.45mm);
  \draw[black!70, line width=0.6pt, dash pattern=on 0.9pt off 0.7pt] (0,0) circle (1.2mm);}
\makeatother

% \Wende{x,y}: Wendepunkt (schwarzer Ring um die Station)
\newcommand\Wende[1]{\draw[black, line width=0.9pt, fill=white] (#1) circle (2.55mm);}

% \St[*]{status}{x}{y}{Platzierung}{Titel}{Detail}
%   *            = Wendepunkt
%   Platzierung  = above | below | left | right | above left | ... (Abstand fest)
% Abstand der Beschriftung zur Marke; Wendepunkte brauchen mehr wegen des Rings
\newcommand\lblsep{2.35mm}
\newcommand\lblsepW{3.3mm}
\NewDocumentCommand\St{s m m m m m m}{%
  \IfBooleanTF{#1}{\Wende{#3,#4}\def\sep{\lblsepW}}{\def\sep{\lblsep}}%
  \Mk{#2}{#3,#4}%
  \node[lbl, #5=\sep] at (#3,#4) {{\bfseries #6}\ifx\relax#7\relax\else\\#7\fi};}

% \Prellbock{x,y}{Winkel}: Ende eines Stumpfgleises (Querbalken)
\newcommand\Prellbock[3][black]{%
  \draw[#1, line width=1.6pt, line cap=butt] ($(#2)+(#3+90:1.7mm)$) -- ($(#2)+(#3-90:1.7mm)$);}

% \Legende: Zeichenerklärung als eigenes Bild (Varianten skalieren mit \scalebox)
\newcommand\Legende{%
\begin{tikzpicture}[x=1cm, y=1cm, every node/.style={font=\fsLbl, anchor=west, inner sep=1pt}]
  \Mk{ok}{0,0}\node at (0.25,0) {Erfolg};
  \Mk{part}{1.45,0}\node at (1.7,0) {Teilerfolg};
  \Mk{fail}{3.25,0}\node at (3.5,0) {Fehlschlag};
  \Mk{old}{5.1,0}\node at (5.35,0) {verworfen / überholt};
  \Wende{7.95,0}\Mk{ok}{7.95,0}\node at (8.25,0) {Wendepunkt};
  \draw[stub, black!70] (10.1,0) -- (10.9,0); \Prellbock[black!70]{10.9,0}{0}
  \node at (11.05,0) {Stumpfgleis = aufgegebener Weg};
  \draw[irrweg, black!70] (15.3,0) -- (16.1,0);
  \node at (16.2,0) {Fehlannahme = Dauer bis zur Widerlegung};
  \node at (20.55,0) {T$n$ = Trainingslauf $n$ · L$n$ = Diagnoselauf $n$ · Zeitachse nicht maßstäblich};
\end{tikzpicture}}

% \BlockerKette: Ursachenkette als Zeile (jede gelöste Ebene legte die nächste frei)
\newcommand\BlockerKette{%
\begin{tikzpicture}[x=1cm, y=1cm,
    bk/.style={draw=black!45, line width=0.5pt, fill=white, rounded corners=1.2mm,
      font=\fsLbl, align=center, text width=31mm, minimum height=9.5mm, inner sep=1.2pt},
    pf/.style={-{Latex[length=1.6mm]}, line width=0.9pt, black!55}]
  \node[font=\fsLbl\bfseries, anchor=east, align=right] at (-0.25,0)
    {Blocker-Kette:\\[-1pt]{\mdseries jede gelöste Ebene}\\[-1pt]{\mdseries legte die nächste frei}};
  \foreach \i/\s/\t/\d in {
      0/ok/{Kameras verdreht}/{behoben, L10–18},
      1/ok/{Geometrie: Finger, Becken, Basis}/{behoben, L38–51},
      2/part/{Würfel 3 cm versetzt}/{korrigiert, Ursache offen, L62},
      3/part/{Bild-Lücke Sim ↔ Echt}/{×13 → halbiert, L66/L83},
      4/part/{Griff}/{2/60 frei gehoben, L97},
      5/open/{Stapeln + Robotertest}/{0/60, offen}} {
    \node[bk] (b\i) at ({\i*3.75+1.75},0) {{\bfseries\t}\\\d};
    \Mk{\s}{{\i*3.75+0.12},0.47}
  }
  \foreach \i [evaluate=\i as \j using int(\i+1)] in {0,...,4} \draw[pf] (b\i.east) -- (b\j.west);
\end{tikzpicture}}

\pagestyle{empty}
\setlength{\parindent}{0pt}

\newcommand\fBody{\fontsize{15}{19}\selectfont\sffamily}
\newcommand\fSmall{\fontsize{12}{15}\selectfont\sffamily}

% \Rahmen{Titel}{Seitenzahl}: Kopfband, Logo, Fußzeile
\newcommand\Rahmen[2]{%
  \coordinate (O) at (current page.north west);
  \fill[lTrain] (O) rectangle ($(O)+(338.67,-19)$);
  \node[anchor=west, text=white, font=\fontsize{24}{28}\selectfont\sffamily\bfseries]
    at ($(O)+(10,-9.5)$) {#1};
  \node[anchor=east, fill=white, rounded corners=1.2mm, inner sep=1.2mm] at ($(O)+(331,-9.5)$)
    {\includegraphics[height=11mm]{../logo_hochschule.jpg}};
  \node[anchor=south west, font=\fontsize{8.5}{10}\selectfont\sffamily, text=black!55]
    at ($(O)+(10,-188)$) {GR00T-Fine-Tuning für Unitree G1 + DEX3 · Berger, Späth, Mücke ·
      Betreuer: Prof. Dr. Matthias Lorenzen · HS Kempten · Stand 01.11.2026};
  \node[anchor=south east, font=\fontsize{8.5}{10}\selectfont\sffamily, text=black!55]
    at ($(O)+(328.67,-188)$) {#2/3};}

\newsavebox\legbox
\newsavebox\icOk \newsavebox\icPart \newsavebox\icFail \newsavebox\icOpen
\begin{document}
\sbox\legbox{\Legende}
\sbox\icOk{\tikz[baseline=-0.7ex]{\Mk{ok}{0,0}}}\sbox\icPart{\tikz[baseline=-0.7ex]{\Mk{part}{0,0}}}
\sbox\icFail{\tikz[baseline=-0.7ex]{\Mk{fail}{0,0}}}\sbox\icOpen{\tikz[baseline=-0.7ex]{\Mk{open}{0,0}}}

% ════════════════════ Folie 1: Gesamtkarte ════════════════════════════
\noindent\begin{tikzpicture}[remember picture, overlay, x=1mm, y=1mm]
  \Rahmen{Projektverlauf auf einen Blick: Mai bis November 2026}{1}
  \node[anchor=north, inner sep=0] at ($(O)+(169.3,-21.5)$)
    {\includegraphics[height=152mm]{build/karte_kurz.pdf}};
  \node[anchor=south, inner sep=0] at ($(O)+(169.3,-184)$) {\scalebox{1.02}{\usebox\legbox}};
\end{tikzpicture}
\newpage

% ════════════════════ Folie 2: Stand heute ════════════════════════════
\noindent\begin{tikzpicture}[remember picture, overlay, x=1mm, y=1mm]
  \Rahmen{Stand heute: erste Griffe, noch kein Stapeln}{2}
  % Diagramm Closed Loop: Anteil Episoden je Stufe
  \begin{scope}[shift={($(O)+(14,-34)$)}]
    \node[anchor=west, font=\fBody\bfseries] at (0,4) {Closed Loop in der Simulation (Szene aus Lauf 70, 120 s je Episode)};
    \def\bw{150}
    \foreach \i/\lab/\a/\na/\b/\nb in {0/erreicht/16/20/58/60, 1/berührt/17/20/60/60,
        2/gehoben/0/20/6/60, 3/frei gehoben/0/20/2/60, 4/gestapelt/0/20/0/60} {
      \pgfmathsetmacro\ya{-10-\i*22}
      \node[anchor=east, font=\fBody] at (40,\ya-3.5) {\lab};
      \fill[black!8] (43,\ya) rectangle ++(\bw,-5.5);
      \fill[black!8] (43,\ya-7) rectangle ++(\bw,-5.5);
      \pgfmathsetmacro\wa{\bw*\a/\na}\pgfmathsetmacro\wb{\bw*\b/\nb}
      \ifnum\a>0 \fill[lTrain!70] (43,\ya) rectangle ++(\wa,-5.5); \fi
      \ifnum\b>0 \fill[lCotr] (43,\ya-7) rectangle ++(\wb,-5.5); \fi
      \node[anchor=west, font=\fSmall] at (\bw+45,\ya-2.75) {\a/\na};
      \node[anchor=west, font=\fSmall\bfseries] at (\bw+45,\ya-9.75) {\b/\nb};
    }
    % Referenz: abgespielte Demo (Orakel, Lauf 69) bei "frei gehoben"
    \draw[black, line width=0.6mm, dash pattern=on 2mm off 1.2mm] ({43+150*0.30},-69) -- ++(0,-20);
    \node[anchor=west, font=\fSmall, align=left] at ({44.5+150*0.30},-71.3)
      {Orakel (abgespielte Demo, Lauf 69): 6/20 = 30\,\%};
    % Legende
    \fill[lTrain!70] (43,-121) rectangle ++(6,-4.5);
    \node[anchor=west, font=\fSmall] at (50,-123.2) {N1.6 vision-v2 @44k, ohne Co-Training (Lauf 82, 20 Lagen)};
    \fill[lCotr] (43,-129) rectangle ++(6,-4.5);
    \node[anchor=west, font=\fSmall] at (50,-131.2) {N1.6 co-trainiert, cube4 @44k (Lauf 97, 60 Lagen)};
  \end{scope}
  % Rechte Spalte: Blocker und nächste Schritte
  \node[anchor=north west, text width=105mm, align=left, font=\fBody] at ($(O)+(222,-27)$) {%
    {\bfseries Was das Co-Training gebracht hat}\par\smallskip
    \fontsize{13.5}{18}\selectfont
    \usebox\icOk\hspace{1.6mm} Fehler auf Sim-Bildern halbiert (0,024 → 0,012)\par
    \usebox\icOk\hspace{1.6mm} Hand öffnet und schließt: 73\,\% der Demo-Spanne\par
    \usebox\icPart\hspace{1.6mm} erste Griffe 2/60 (noch nicht signifikant)\par
    \usebox\icPart\hspace{1.6mm} Echtbild-Fehler N1.6 +50–60\,\%\par
    \bigskip
    {\fBody\bfseries Was noch fehlt}\par\smallskip
    \usebox\icFail\hspace{1.6mm} gestapelt: in keinem Lauf\par
    \usebox\icFail\hspace{1.6mm} kein Test am echten Roboter\par
    \usebox\icOpen\hspace{1.6mm} Ursache des 3-cm-Versatzes offen\par
    \bigskip
    {\fBody\bfseries Nächste Schritte}\par\smallskip
    1.\ Test am G1: cube4 gegen vision-v2\par
    2.\ mehr Render-Episoden statt höherer Mischung\par
    3.\ RL auf dem co-trainierten Modell};
\end{tikzpicture}
\newpage

% ════════════════════ Folie 3: Irrwege und Lessons Learned ═══════════
\noindent\begin{tikzpicture}[remember picture, overlay, x=1mm, y=1mm]
  \Rahmen{Wie lange galt ein Irrtum? Und was wir daraus lernen}{3}
  \begin{scope}[shift={($(O)+(14,-32)$)}]
    % Zeitachse: 1 Tag = 0,9 mm, Tag 0 = 01.05.2026
    \foreach \d/\m in {0/Mai,31/Jun,61/Jul,92/Aug,123/Sep,153/Okt,184/Nov} {
      \draw[black!25, line width=0.3mm] ({58+\d*0.78},2) -- ({58+\d*0.78},-112);
      \node[anchor=south, font=\fSmall, text=black!60] at ({58+\d*0.78},2) {\m};
    }
    \foreach \i/\lab/\a/\b/\c/\txt in {
        0/{Kameras verdreht}/29/99/lSim/{70 Tage, L10–18},
        1/{Split nicht aktiv}/32/104/lTrain/{72 Tage, KISSKI},
        2/{Finger-Vorzeichen}/34/115/lSim/{≥ 81 Tage, L38–46},
        3/{Demo nie als Referenz}/33/166/lDiag/{133 Tage, L57/62},
        4/{SigLIP als Gap-Maß}/34/166/lDiag/{132 Tage, L65/66},
        5/{U-Kurve @30\,000}/105/164/lTrain/{59 Tage, 20×750}} {
      \pgfmathsetmacro\yy{-10-\i*19}
      \node[anchor=east, font=\fSmall\bfseries] at (55,\yy) {\lab};
      \fill[\c!85, rounded corners=1.8mm] ({58+\a*0.78},\yy+4) rectangle ({58+\b*0.78},\yy-4);
      \node[anchor=west, font=\fSmall] at ({59.5+\b*0.78},\yy) {\txt};
    }
    \node[anchor=north west, font=\fSmall, text=black!60, text width=200mm] at (0,-116)
      {Balken: vom ersten Auftreten bis zur Widerlegung. Jeder Irrtum lieferte plausible Zahlen ohne
       Fehlermeldung.};
  \end{scope}
  \node[anchor=north west, text width=82mm, align=left, font=\fontsize{13.5}{17}\selectfont\sffamily] at ($(O)+(246,-28)$) {%
    {\fBody\bfseries Lessons Learned}\par\medskip
    \textbf{Referenz zuerst:} die abgespielte Demo muss in der Sim greifen.\par\medskip
    \textbf{Wirkung messen, nicht Ähnlichkeit:} Teacher-Forcing auf Sim-Bildern statt SigLIP.\par\medskip
    \textbf{Szene unabhängig prüfen:} gegen Echtbilder, nicht mit demselben Code.\par\medskip
    \textbf{Genug Episoden, feste Seeds}, Regeln vorab als Zahl.\par\medskip
    \textbf{Laufenden Prozess prüfen:} ein Schalter ist kein Beleg.\par\medskip
    \textbf{Laborbuch mit Laufnummern} führen.};
\end{tikzpicture}
\end{document}
